% 原书图 2.23 的定性重绘：孤立的 N 个阱与一个共同势阱。
\begin{tikzpicture}[x=0.85cm,y=0.6cm,>=Stealth,font=\small]
\foreach \x in {0,1.5,3,6.5}{
  \draw[thick] (\x,2.1)--(\x,0)--(\x+0.95,0)--(\x+0.95,2.1);
  \draw (\x+0.48,0.08) circle (0.065);
}
\node at (5.15,0.8) {$\cdots$};
\draw[|<->|] (0,-0.4)--(0.95,-0.4) node[midway,below] {$a$};
\node at (3.75,-0.65) {（a）};
\begin{scope}[yshift=-3.2cm]
\draw[thick] (0,2.1)--(0,0)--(7.45,0)--(7.45,2.1);
\foreach \x in {0.3,1.05,1.8,2.55,3.3,5.2,5.95,6.7}{
  \draw (\x,0.08) circle (0.065);
}
\node at (4.2,0.12) {$\cdots$};
\draw[|<->|] (0,-0.4)--(7.45,-0.4) node[midway,below] {$Na$};
\node at (3.75,-1.2) {（b）};
\end{scope}
\end{tikzpicture}
